% 平坦 vs 尖锐极小（泛化角度）：双子图
% 左：平坦（宽碗 + 训练/测试损失接近）
% 右：尖锐（窄碗 + 测试损失远高于训练）
\documentclass[tikz,border=4pt]{standalone}
\usepackage{ctex}
\usepackage{amsmath}
\usetikzlibrary{calc, arrows.meta, patterns, backgrounds}

\begin{document}
\begin{tikzpicture}[scale=1.0, thick]

% ─── 左子图：平坦极小 ────────────────────────────────
\begin{scope}[xshift=0cm]
  % 背景边框
  \draw[black!20, thin, rounded corners=3pt] (-0.5, -0.8) rectangle (5.5, 4.0);
  \node[font=\small\bfseries, align=center] at (2.5, 3.7) {平坦极小（Flat Minimum）};

  % 坐标轴
  \draw[->, gray] (-0.2, 0) -- (5.2, 0) node[right, font=\footnotesize] {$\theta$};
  \draw[->, gray] (0, -0.2) -- (0, 3.5) node[above, font=\footnotesize] {$\mathcal{L}$};

  % 训练损失曲线（实线，蓝色）：宽碗形
  \draw[blue, very thick]
    plot[smooth, tension=0.7]
    coordinates {(0.3, 3.2) (0.9, 2.0) (1.6, 0.8) (2.5, 0.2) (3.4, 0.8) (4.1, 2.0) (4.7, 3.2)};
  \filldraw[blue] (2.5, 0.2) circle (2.5pt);

  % 测试损失曲线（虚线，蓝色较浅）：形状类似，略高
  \draw[blue!60, very thick, dashed]
    plot[smooth, tension=0.7]
    coordinates {(0.3, 3.4) (0.9, 2.2) (1.6, 1.0) (2.5, 0.45) (3.4, 1.0) (4.1, 2.2) (4.7, 3.4)};
  \filldraw[blue!60] (2.5, 0.45) circle (2.5pt);

  % 极小点标注
  \draw[gray, thin] (2.5, 0) -- (2.5, 0.2);
  \node[below, font=\footnotesize] at (2.5, -0.05) {$\theta^*$};

  % 泛化差距很小的标注（双向箭头）
  \draw[<->, black!50, thin] (2.75, 0.2) -- (2.75, 0.45);
  \node[right, font=\footnotesize, black!50] at (2.8, 0.32) {$\Delta$ 小};

  % 宽度标注
  \draw[blue!40, dashed, thin] (1.2, 0) -- (1.2, 0.7);
  \draw[blue!40, dashed, thin] (3.8, 0) -- (3.8, 0.7);
  \draw[blue!40, <->, thin] (1.2, 0.55) -- (3.8, 0.55) node[midway, above, font=\footnotesize, blue!70] {宽平区域（$r$ 大）};

  % 分布偏移示意
  \draw[->, gray!60, thin, dashed] (2.5, 0.2) -- (3.2, 0.2);
  \node[right, font=\tiny, gray!60] at (3.2, 0.2) {分布偏移};

  % 图例
  \draw[blue, very thick] (0.3, 3.0) -- (1.1, 3.0);
  \node[right, font=\footnotesize, blue] at (1.1, 3.0) {训练损失};
  \draw[blue!60, very thick, dashed] (0.3, 2.65) -- (1.1, 2.65);
  \node[right, font=\footnotesize, blue!60] at (1.1, 2.65) {测试损失};

  \node[font=\footnotesize, blue, align=center] at (2.5, -0.5) {$\mathcal{L}_\text{train} \approx \mathcal{L}_\text{test}$};
\end{scope}

% ─── 右子图：尖锐极小 ───────────────────────────────
\begin{scope}[xshift=6.5cm]
  % 背景边框
  \draw[black!20, thin, rounded corners=3pt] (-0.5, -0.8) rectangle (5.5, 4.0);
  \node[font=\small\bfseries, align=center] at (2.5, 3.7) {尖锐极小（Sharp Minimum）};

  % 坐标轴
  \draw[->, gray] (-0.2, 0) -- (5.2, 0) node[right, font=\footnotesize] {$\theta$};
  \draw[->, gray] (0, -0.2) -- (0, 3.5) node[above, font=\footnotesize] {$\mathcal{L}$};

  % 训练损失曲线（实线，红色）：窄尖碗形
  \draw[red, very thick]
    plot[smooth, tension=1.2]
    coordinates {(1.5, 3.2) (1.8, 2.2) (2.1, 1.0) (2.5, 0.2) (2.9, 1.0) (3.2, 2.2) (3.5, 3.2)};
  \filldraw[red] (2.5, 0.2) circle (2.5pt);

  % 测试损失曲线（虚线，红色较深）：更高，且受分布偏移影响更大
  \draw[red!80, very thick, dashed]
    plot[smooth, tension=1.1]
    coordinates {(1.3, 3.2) (1.7, 2.4) (2.0, 1.4) (2.5, 0.9) (3.0, 1.4) (3.3, 2.4) (3.7, 3.2)};
  \filldraw[red!80] (2.5, 0.9) circle (2.5pt);

  % 极小点标注
  \draw[gray, thin] (2.5, 0) -- (2.5, 0.2);
  \node[below, font=\footnotesize] at (2.5, -0.05) {$\theta^*$};

  % 泛化差距很大的标注
  \draw[<->, black!50, thin] (2.75, 0.2) -- (2.75, 0.9);
  \node[right, font=\footnotesize, black!50] at (2.8, 0.55) {$\Delta$ 大！};

  % 宽度标注（很窄）
  \draw[red!40, dashed, thin] (2.1, 0) -- (2.1, 0.85);
  \draw[red!40, dashed, thin] (2.9, 0) -- (2.9, 0.85);
  \draw[red!40, <->, thin] (2.1, 0.7) -- (2.9, 0.7) node[midway, above, font=\footnotesize, red!60] {窄（$r$ 小）};

  % 分布偏移后损失急增示意
  \draw[->, red!50, thin] (2.5, 0.2) -- (2.8, 0.2);
  \node[right, font=\tiny, red!50] at (2.8, 0.2) {偏移};
  \draw[->, red!80, thin] (2.8, 0.2) -- (2.8, 0.85);
  \node[right, font=\tiny, red!80] at (2.85, 0.55) {急增};

  % 图例
  \draw[red, very thick] (0.3, 3.0) -- (1.1, 3.0);
  \node[right, font=\footnotesize, red] at (1.1, 3.0) {训练损失};
  \draw[red!80, very thick, dashed] (0.3, 2.65) -- (1.1, 2.65);
  \node[right, font=\footnotesize, red!80] at (1.1, 2.65) {测试损失};

  \node[font=\footnotesize, red, align=center] at (2.5, -0.5) {$\mathcal{L}_\text{test} \gg \mathcal{L}_\text{train}$};
\end{scope}

\end{tikzpicture}
\end{document}
